\documentclass[11pt]{article}
\usepackage[margin=1in]{geometry}
\usepackage{amsmath,amssymb,amsthm,booktabs}
\usepackage[hidelinks]{hyperref}
\setlength{\emergencystretch}{2em}
\newtheorem{theorem}{Theorem}
\newtheorem{proposition}[theorem]{Proposition}
\newtheorem{corollary}[theorem]{Corollary}
\newcommand{\F}{\mathbb F}
\newcommand{\mr}{\operatorname{mr}}
\newcommand{\rank}{\operatorname{rank}}
\title{Singular affine certificates and exact $GL_3(\F_2)$ quotient tests}
\author{Internal research working proof}
\date{16 November 2026}
\begin{document}
\maketitle

There is a concrete representation architecture missed by the pinned
external four-dimensional search: homogeneous affine extensions can
have dependent meridian image directions and still have no common
fixed vector. An explicit example gives a sharp certificate for every
$T(4,7(2r+1))$, $r\ge0$. These are known torus-knot MRC cases
\cite{rz}; no new knot equality or priority claim is made.

This observation motivated a frozen, complete $GL_3(\F_2)$ quotient
test on all 293 names in the pinned external open list. Independent
replay finds 191 quotient kernels on 154 knots. Neither the binary
affine lifts nor the tested characteristic-seven reflection lifts
give a new sharp case. The original sixteen targets are retained as
a fixed prefix. The full classical MRC goal remains open.

\section{The homogeneous affine mechanism}

Let $\Gamma$ be a knot group, $\rho:\Gamma\to GL(V)$ a representation
over a field $k$, and $c\in Z^1(\Gamma,V)$, with convention
$c(gh)=c(g)+\rho(g)c(h)$. Homogenization gives
\[
 \widehat\rho(g)=\begin{pmatrix}\rho(g)&c(g)\\0&1\end{pmatrix}
 \quad\text{on }V\oplus k.
\]

\begin{proposition}\label{prop:affine}
Suppose $V^\Gamma=0$, $c$ is nonprincipal, and for a meridian $m$,
$\rank(\rho(m)-I)=1$ and $c(m)\in\operatorname{im}(\rho(m)-I)$.
Then every meridian of $\widehat\rho$ has rank-one difference from
the identity, its common fixed space is zero, and
$\mr(K)\ge\dim V+1$.
\end{proposition}
\begin{proof}
The image of $\widehat\rho(m)-I$ is the original one-dimensional
image in $V\oplus0$. Conjugacy gives the assertion for every meridian.
A fixed vector $(v,0)$ has $v=0$. A fixed vector $(v,t)$ with
$t\ne0$ would give $c(g)=t^{-1}(v-\rho(g)v)$ for every $g$,
making $c$ principal. Thus the common fixed space is zero.
The fixed hyperplanes of any $s$ generating meridians have
intersection of codimension at most $s$; their intersection is the
global fixed space. Hence $s\ge\dim V+1$.
\end{proof}

This is a direct application of the retained admissible-cocycle and
rank-one fixed-space bounds, without a common bilinear-form
hypothesis. In dimension three a nonprincipal admissible class
gives a four-dimensional certificate even though all homogeneous
meridian image lines lie in the same three-space.

\section{An explicit sharp control and a uniform known family}

Over $\F_2$ put
\[
 X=\begin{pmatrix}1&0&0&1\\1&1&0&0\\0&1&1&0\\0&0&0&1\end{pmatrix},
 \qquad
 Y=\begin{pmatrix}0&1&0&0\\0&0&1&0\\1&1&0&0\\0&0&0&1\end{pmatrix}.
\]
Direct multiplication gives
\[
 X^4=Y^7=I,\qquad
 M=X^{-1}Y^2=
 \begin{pmatrix}0&0&1&1\\1&1&1&1\\1&0&0&1\\0&0&0&1\end{pmatrix}.
\]
The three nonzero rows of $M-I$ coincide, so its rank is one.
The equations $(X-I)v=(Y-I)v=0$ force all four coordinates of $v$
to vanish: the first matrix forces $v_4=v_1=v_2=0$, and the second
then forces $v_3=0$. This already proves the needed lower bound;
image enumeration is an additional control.

\begin{theorem}\label{thm:torus}
For every positive odd multiple $q$ of seven, the above matrices
give a sharp rank-one certificate for $T(4,q)$, and
$\mr(T(4,q))=b(T(4,q))=4$.
\end{theorem}
\begin{proof}
Use $\Gamma=\langle x,y\mid x^4=y^q\rangle$ with preferred
meridian $x^a y^b$, where $a q+4b=1$, in the standard conventions
of \cite[Section 5.1]{cornwell}. Since $q$ is odd, $a$ is odd;
reduction modulo seven gives $b\equiv2\pmod7$.
Choose $\epsilon\in\{1,-1\}$ with
$\epsilon a\equiv-1\pmod4$. Send $x$ to $X^\epsilon$ and $y$ to
$Y$. The relation holds, and the meridian maps to $M$.
The common fixed space is zero, also with $X$ replaced by its
inverse. The fixed-hyperplane argument gives $\mr\ge4$.
The standard four-strand torus braid gives $b\le4$, and
$\mr\le b$ completes the proof.
\end{proof}

For $T(4,7)$ the underlying linear three-dimensional image is
$GL_3(\F_2)$, and exact affine enumeration gives the split group
$AGL_3(\F_2)$ of order 1344, transitive on eight points. The
admissible cocycle space has dimension four, its principal subspace
has dimension three, and the displayed translation is nonprincipal.
These facts have separate dense-matrix and eight-point permutation
replay. They are not a claim about the distinct nonsplit group
$2^3.L_3(2)$ listed by ATLAS.

The natural $GL_3(\F_2)$ action preserves no nonzero alternating
form: the radical of an invariant alternating form is invariant,
whereas an alternating form in odd dimension is singular and the
natural action is irreducible. This affine certificate therefore
uses a broader mechanism than the original common-form criterion.
More strongly, when $3\nmid q$, that criterion has no realization
on $T(4,q)$: $b\le4$ forces its dimension $k\le3$, a nonzero
alternating form then has rank two, and the retained rank-two
theorem forces an $S_3$ transposition quotient. Such a quotient
would give nontrivial Fox three-colouring, contrary to
$\det T(4,q)=q$. These are comparisons of certificate classes;
Rost--Zieschang already proved torus-knot MRC \cite{rz}.

\section{The exact external search limitation}

At the pinned commit listed in the bibliography, the external
program \texttt{search/glsearch.py} fixes four seed
directions to $e_1,\ldots,e_4$ and explicitly discards the seed
functional matrix $A$ when $\det A=0$ \cite{search}.
This searches a restricted class of representations.

Indeed, write four homogeneous meridians as $I+u_i f_i$, let
$U$ have columns $u_i$, and let $F$ have rows $f_i$. Every image
line of an affine extension lies in $V\oplus0$. For $\dim V=3$,
\[
 \rank U\le3,\qquad \rank(FU)\le3.
\]
Thus no choice of four meridians, scalings or basis change can
produce the search's independent directions or invertible pairing
matrix. Zero common fixed space instead concerns the functionals:
it does not require the image directions to be independent.
The $T(4,7)$ replay exhibits four conjugate meridians with
independent functionals but pairing rank three.

This is a search-coverage limitation, not an error in the valid
rank-one criterion or in the 243 deposited certificates. The
three-seed program \texttt{search/w3.py} performs bounded random
diagram moves and does not certify absence of a three-bridge
presentation. In the archived KnotInfo inputs the 293 open names
have bridge interval $[3,4]$; their bridge number is not assumed
to be four. No current live-list or exhaustive-geometric-search
claim is made.

\section{A complete frozen \texorpdfstring{$GL_3(\F_2)$}{GL3(F2)} quotient panel}

Let $G=\pi_1(\Sigma_2(K))$. Retained four-seed diagram certificates
give three generators for $G$. Enumerating generating triples of
$Q=GL_3(\F_2)\cong PSL_2(\F_7)$ modulo its 336 automorphisms
therefore exhausts epimorphisms up to quotient kernel. There are
13,368 such triple orbits. The independently counted pair capacity
is 57, so 58 distinct kernels on a single $G$ would force rank
at least three. This is the same classical simple-quotient capacity
principle as the user's $A_5$ threshold, with no pooling across
knots.

All 293 pinned names, including the original sixteen, were frozen
before the enlarged panel ran. The exact results are
\[
\begin{array}{c|rrrrr}
\text{kernels per knot}&0&1&2&3&7\\\hline
\text{number of knots}&139&123&29&1&1
\end{array}
\]
The maximum seven occurs at \texttt{13n\_4639}; three occur at
\texttt{13n\_4645}. There are 191 kernels on 154 knots, far below
the same-source threshold. The deck action is inner for 29 maps
and outer for 162. Both natural and dual binary ordinary cocycle
spaces have dimension four on 179 maps and five on twelve.
The eight-dimensional Steinberg space has dimension eight on
165 maps and ten on 26.

\begin{corollary}\label{cor:modules}
For every retained epimorphism $G\twoheadrightarrow Q$, every
field $k$ of characteristic two and every finite-dimensional
$kQ$-module $V$,
\[
 \dim Z^1(G,V)\le\tfrac53\dim V.
\]
Consequently the ordinary linear-cocycle bound cannot force
$\rank G\ge3$ through any of these epimorphisms.
\end{corollary}
\begin{proof}
The trivial-module cocycle space is zero in characteristic two,
as independently checked on each retained presentation. The other
absolute simples have dimensions $3,3,8$ \cite{atlas}. The displayed
values bound their normalized cocycle dimensions by $5/3$.
For $0\to U\to V\to W\to0$, the map on cocycles has kernel
$Z^1(G,U)$, so
$\dim Z^1(G,V)\le\dim Z^1(G,U)+\dim Z^1(G,W)$.
Composition series and scalar extension give the assertion even
for nonsplit extensions and arbitrary coefficient fields of
characteristic two.
\end{proof}

The simple-model algebra spans have dimensions nine and 64,
proving absolute simplicity; duality handles the second natural
module. Independently enumerated two-regular conjugacy classes
have orders/sizes $(1,1),(3,56),(7,24),(7,24)$. The standard
Brauer class-count theorem supplies completeness of these four
absolute simple models, also consistent with the primary ATLAS
list. The classification theorem remains an explicit import.
The corollary covers full $Q$ images, not all proper-image
representations into $GL_3(\F_2)$ on this enlarged panel, and not
other characteristics or finite groups. It gives no rank-two
upper certificate for $G$.

For inner deck action, the natural and dual modules extend to
binary rank-one meridian bases. All 58 associated root-cocycle
spaces have dimension three. For both deck types, identify $Q$
with $PSL_2(\F_7)$ and send a projective meridian $A$ to
$-\operatorname{Ad}(A)$ on traceless $2\times2$ matrices. Each is
an involution with rank-one difference from identity. Outer deck
action gives the determinant-character-twisted $PGL_2(\F_7)$
action; inner deck action uses the separate negative sign and
gives $C_2\times PSL_2(\F_7)$. All 191 root spaces again have
dimension three. These are complete tests of these specified
bases and fields, not all characteristic-seven coefficient modules.

\section{Replay, provenance and remaining work}

The retained plans, catalogues, matrices and kernel bases are
sealed by \texttt{gl3f2\_verification.json}. The verifier constructs
an independent M\"obius-permutation $PSL_2(\F_7)$ model, verifies
all products under an explicit identification, marks full tuple
orbits, and uses iterative crossing closure to reproduce every
accepted map on all 293 diagrams. Ordinary matrices are rebuilt
by reverse affine-point evaluation. The Steinberg model from
Fano-building cycles is transported to traceless endomorphisms
by an explicit invertible intertwiner; all 168 matrices agree.
Independent conjugation constructs the characteristic-seven
adjoint action and checks every root matrix and nullspace.
The positive torus control also has exact dense/permutation replay.
The source PD parser and retained diagram conventions are shared;
this is independent algebraic replay, not independent knot
identification or a formal proof assistant verification.

All 191 maps, including the first-sixteen prefix, fail the sharp
lift gates tested here. Known-family controls do not close a new
open case. A next quotient or representation test needs a new
structural reason for a rank gain. Proper images, other finite
targets, nonlinear group-rank obstructions and genuinely different
meridian components remain outside these exclusions. Actual cover
ranks, target MRC equalities and full MRC remain unresolved.
The uniform working proofs are unformalized and externally
unreviewed.

\begin{thebibliography}{9}
\bibitem{rz} Rost--Zieschang, Meridional generators and plat
presentations of torus links, J. London Math. Soc. 35 (1987),
551--562. \url{https://doi.org/10.1112/jlms/s2-35.3.551}.
\bibitem{cornwell} Cornwell, Knot contact homology and representations
of knot groups, Section 5.1.
\url{https://arxiv.org/abs/1303.4943v5}.
\bibitem{atlas} ATLAS of Finite Group Representations,
$L_2(7)=L_3(2)$, representation and automorphism lists.
\url{https://brauer.maths.qmul.ac.uk/Atlas/lin/L32/}.
\bibitem{search} Bhattacharjee, unrefereed certificate project,
pinned search programs and deposited certificates.
\url{https://github.com/anaxia-research/conjecture-to-choose/blob/06cbbb7e2737e1a11085b64e0ead021d63c550c3/search/glsearch.py}.
\url{https://zenodo.org/records/23561157}.
\end{thebibliography}
\end{document}
